%=====================================================================
\chapter{Bayesian Networks}\label{ch:bayesnets}
%=====================================================================
\locator{5}{Part V --- Reasoning Under Uncertainty}

\begin{outcomes}
By the end of this chapter you will be able to:
\begin{tight}
  \item \textbf{Read} a Bayesian network as a set of conditional independence
        claims, and \textbf{count} the numbers it needs against the full joint.
  \item \textbf{Compute} a posterior by enumeration and by variable elimination.
  \item \textbf{Estimate} a posterior by rejection sampling and by likelihood
        weighting, and \textbf{say} when each is preferable.
  \item \textbf{Recognise} \emph{explaining away}, and \textbf{explain} why it
        cannot be produced by any rule-based system.
  \item \textbf{Judge} when variable elimination is worth its overhead.
\end{tight}
\end{outcomes}

\begin{prereq}
\chref{uncertainty} --- Bayes' rule, marginalisation, and above all conditional
independence, which is the single idea this whole representation is built on.
\end{prereq}

\begin{keyterms}
Bayesian network $\bullet$ directed acyclic graph $\bullet$ conditional probability
table $\bullet$ parent $\bullet$ chain rule $\bullet$ enumeration $\bullet$ factor
$\bullet$ variable elimination $\bullet$ rejection sampling $\bullet$ likelihood
weighting $\bullet$ explaining away $\bullet$ diagnostic and causal direction
\end{keyterms}

\section{Diagnosing a Defect}

A release has been rolled back. Why?

The team believes three things make defects likely: a rushed specification, a
junior developer working unsupervised, and weak test coverage. A defect in turn
makes two things likely: an emergency hotfix, and a rollback. Nobody observes the
defect directly --- they observe the rollback and want to know what caused it.

That is a diagnosis problem, and its structure is what a \term{Bayesian network}
represents.

\begin{definitionbox}[Bayesian network]
A \term{Bayesian network} is a directed acyclic graph in which each node is a random
variable and each node carries a \term{conditional probability table} giving
$P(\text{node} \mid \text{its parents})$.

The graph asserts that \emph{each variable is conditionally independent of its
non-descendants given its parents}. That assertion licenses the chain rule
\[
  P(x_{1}, \ldots, x_{n}) \;=\; \prod_{i=1}^{n}
  P\bigl(x_{i} \mid \mathrm{parents}(x_{i})\bigr)
\]
so the full joint distribution is recoverable from the tables --- which is why the
tables are all you have to store, and all anybody has to estimate.
\end{definitionbox}

\dsfig{%
\begin{tikzpicture}[font=\scriptsize,
  n/.style={circle, draw=#1, line width=1.1pt, fill=#1!12, text=#1!55!black,
            font=\scriptsize\bfseries, minimum size=1.35cm, align=center,
            inner sep=1pt},
  e/.style={-{Stealth[length=2mm]}, primary!60, line width=1pt},
  cpt/.style={rounded corners=2pt, draw=#1!50, fill=#1!5, text=#1!55!black,
              font=\tiny, align=flush left, inner sep=4pt}]

\node[n=goldhl]  (r) at (-3.4,2.0) {Rushed\\spec};
\node[n=goldhl]  (j) at (0,2.0)    {Junior\\dev};
\node[n=goldhl]  (w) at (3.4,2.0)  {Weak\\tests};
\node[n=accent]  (d) at (0,0)      {Defect};
\node[n=greenhl] (h) at (-2.2,-2.0){Hotfix};
\node[n=greenhl] (b) at (2.2,-2.0) {Rollback};

\foreach \a in {r,j,w}{\draw[e] (\a) -- (d);}
\draw[e] (d) -- (h);
\draw[e] (d) -- (b);

\node[cpt=goldhl, anchor=east] at (-4.4,2.0)
  {$P(\mathit{Rushed}) = 0.30$\\
   $P(\mathit{Junior}) = 0.40$\\
   $P(\mathit{Weak}) = 0.25$};

\node[cpt=accent, anchor=west] at (1.5,0.1)
  {$P(\mathit{Defect} \mid R,J,W)$\\[1pt]
   \texttt{T T T} \ 0.95 \quad \texttt{F T T} \ 0.70\\
   \texttt{T T F} \ 0.80 \quad \texttt{F T F} \ 0.40\\
   \texttt{T F T} \ 0.75 \quad \texttt{F F T} \ 0.30\\
   \texttt{T F F} \ 0.50 \quad \texttt{F F F} \ 0.05};

\node[cpt=greenhl, anchor=east] at (-3.3,-2.0)
  {$P(H \mid D)$\\ \texttt{T} \ 0.90\\ \texttt{F} \ 0.10};
\node[cpt=greenhl, anchor=west] at (3.3,-2.0)
  {$P(B \mid D)$\\ \texttt{T} \ 0.70\\ \texttt{F} \ 0.05};

\node[rounded corners=3pt, fill=lightbg, draw=primary!30, line width=0.9pt,
      text=primary!70!black, font=\scriptsize, text width=6.0cm,
      align=flush left, inner sep=6pt, anchor=north] at (0,-3.2)
  {\textbf{The full joint over six binary variables needs $2^{6} = 64$ numbers.
   This network needs 15} --- and each of the 15 is a quantity somebody could
   actually estimate from records, which the 64 are not.};
\end{tikzpicture}}%
{The diagnosis network. Causes point to effects, so the arrows run the way the
world does; inference then runs against them, from rollback back to
cause.}{bayesnet}

\begin{alertbox}[What the missing arrows say]
The content of a Bayesian network is in the edges it \emph{lacks}.

There is no arrow from \emph{Hotfix} to \emph{Rollback}, which asserts that once you
know whether there is a defect, learning that a hotfix was issued tells you nothing
further about whether a rollback happened. There is no arrow between \emph{Rushed}
and \emph{Junior}, asserting they are independent before any evidence.

Those are substantive claims about the domain and they can be wrong. If rushed
specifications are usually handed to junior developers, the second claim is false
and the network will give confidently wrong answers. \emph{Drawing the graph is the
modelling; the numbers are the easy part.}
\end{alertbox}

\section{Exact Inference by Enumeration}

$P(\mathit{Rushed} \mid \mathit{Rollback})$ follows directly from the definition:
\[
  P(r \mid b) \;=\; \alpha \sum_{j,w,d,h} P(r)P(j)P(w)
  P(d \mid r,j,w) P(h \mid d) P(b \mid d)
\]
Sum the joint over every hidden variable, for each value of the query, and
normalise. On this network that is $32$ terms, and the answer is
\[
  P(\mathit{Rushed} \mid \mathit{Rollback}) \;=\; 0.4899
\]
against a prior of $0.30$. A rollback makes a rushed specification substantially
more likely --- not because any rule says so, but because the arithmetic of the
graph says so.

Enumeration is exact and it is $O(2^{n})$. That is \chref{propositional}'s model
checking again, and it fails for the same reason.

\section{Variable Elimination}

Enumeration recomputes the same subexpressions over and over. \term{Variable
elimination} computes each once, by pushing sums inward and caching the results as
\term{factors}.

A factor is a table over some subset of the variables. The procedure is: build one
factor per node; then for each hidden variable, multiply together every factor
mentioning it and sum it out, leaving a new factor; finally multiply what remains
and normalise.

\begin{alertbox}[On this network, variable elimination is slower]
It uses $74$ operations where enumeration uses $32$ terms, and returns the same
$0.4899$. The bookkeeping costs more than the sharing saves, because the network has
six nodes and there is almost nothing to share.

This is the same lesson as \chref{csp}'s AC-3 and \chref{planning}'s heuristic, and
it is worth stating as a rule: \emph{an algorithm whose advantage is asymptotic has
no advantage at the size where you first test it}. Measure at the size you will
deploy at.
\end{alertbox}

Deploy it at scale and the picture inverts completely. Take a chain of services,
each of which may be delayed by the one before it, and ask about the first given
evidence about the last:

\dsfig{%
\begin{tikzpicture}
\begin{axis}[
  width=11.8cm, height=5.6cm,
  xlabel={services in the chain, $n$},
  ylabel={arithmetic operations},
  xlabel style={font=\scriptsize}, ylabel style={font=\scriptsize},
  tick label style={font=\tiny},
  ymode=log, log basis y=10,
  xmin=5, xmax=19, ymin=10, ymax=2e4,
  grid=both, major grid style={gray!18}, minor grid style={gray!8},
  legend style={font=\tiny, at={(0.03,0.97)}, anchor=north west,
                draw=gray!40, fill=white},
  legend cell align=left]

\addplot[accent, line width=1.3pt, mark=*, mark size=1.6pt]
  coordinates {(6,32) (10,512) (14,8192)};
\addlegendentry{enumeration --- doubles with every service}

\addplot[greenhl, line width=1.3pt, mark=square*, mark size=1.6pt]
  coordinates {(6,58) (10,122) (14,186) (18,250)};
\addlegendentry{variable elimination --- linear}

\node[font=\tiny\itshape, text=accent, anchor=west] at (axis cs:14.4,8192)
  {enumeration is already infeasible};
\node[font=\tiny\itshape, text=greenhl!50!black, anchor=west]
  at (axis cs:14.6,250) {250 operations at $n=18$};
\end{axis}
\end{tikzpicture}}%
{Enumeration against variable elimination on a chain. At six nodes enumeration
wins; at fourteen it needs $8{,}192$ terms against $186$ operations, and beyond
that it stops being an option.}{ve-scaling}

Variable elimination is not always polynomial --- its cost is governed by the
largest intermediate factor, which depends on the graph's structure and on the
elimination order. On a chain it is linear; on a densely connected graph it is
exponential again, and exact inference in general Bayesian networks is NP-hard.

\section{Approximate Inference by Sampling}

When exact inference is out of reach, generate samples from the network and count.

\term{Rejection sampling} is the obvious method: sample every variable in
topological order using its conditional probability table, discard any sample that
disagrees with the evidence, and take the proportion among the survivors.

It is correct, and it wastes almost everything. Here the evidence
$\mathit{Rollback} = \text{T}$ has probability about $0.30$, so roughly $70\%$ of
samples are thrown away: $20{,}000$ samples leave about $6{,}000$ usable. With
rarer evidence --- and diagnostic evidence is usually rare, since that is why it is
interesting --- the rejection rate approaches $100\%$ and the method stops working.

\term{Likelihood weighting} fixes this. Do not sample the evidence variables at all:
fix them to their observed values, sample only the rest, and weight each sample by
the probability of the evidence given what was sampled. Every sample now counts.

\smallskip
\noindent\begin{longtable}{@{}L{4.4cm}L{2.7cm}L{2.7cm}L{4.4cm}@{}}
\toprule
\rowcolor{primary!15}
\textbf{Method} & \textbf{Answer} & \textbf{Error} & \textbf{Note} \\
\midrule
\endhead
Enumeration (exact) & 0.48990 & --- & 32 terms \\
Variable elimination (exact) & 0.48990 & --- & 74 operations \\
Rejection, 20{,}000 samples & 0.4952 & 0.0053 & $6{,}034$ kept; $70\%$ discarded \\
Likelihood weighting, 20{,}000 & 0.4913 & \textbf{0.0014} & every sample used \\
\bottomrule
\caption{$P(\mathit{Rushed} \mid \mathit{Rollback})$ four ways. The two exact
methods agree to every digit; the two approximate ones differ mainly in how much
of their work they throw away.}\label{tab:bn-inference}
\end{longtable}

\section{Explaining Away}

Here is the behaviour that justifies the whole apparatus, and that no rule base of
\chref{expert} can produce.

Learn that there \emph{is} a defect. Belief in a rushed specification rises from
$0.30$ to $0.5281$ --- the defect needs explaining, and a rushed specification is
one of the explanations.

Now additionally learn that a junior developer was working unsupervised. Belief in a
rushed specification \emph{falls}, to $0.4304$. The junior developer already
accounts for the defect, so less of the explanatory work is left for the
specification to do.

\begin{conceptbox}
Two causes of a common effect are independent \emph{until} the effect is observed,
and become dependent once it is --- negatively so. That is \term{explaining away},
and it is the reason to build a network rather than a rule base.

Notice what a production rule would have to say. \textsc{if} defect \textsc{and}
junior developer \textsc{then} rushed specification is \emph{less} likely --- a rule
whose conclusion is weakened by adding a condition, which no forward-chaining engine
can express, and which \chref{expert}'s specificity strategy would actively get
wrong. The behaviour falls out of the graph and the arithmetic without anybody
writing it down, and it would have to be anticipated and hand-coded in any
rule-based system.

It is also the formal version of something every incident reviewer does: on finding
one sufficient cause, they become less interested in the others. The network says
exactly how much less.
\end{conceptbox}

\begin{worked}[Computing $P(\mathit{Rushed} \mid \mathit{Rollback})$ by hand]
The query has one evidence variable, and $\mathit{Junior}$, $\mathit{Weak}$,
$\mathit{Defect}$ and $\mathit{Hotfix}$ are hidden.

\smallskip
\textbf{Step 1: drop \emph{Hotfix}.} It is not in the query and not in the
evidence, and nothing else depends on it. Summing it out gives
$\sum_{h} P(h \mid d) = 1$ for either value of $d$, so it contributes a factor of
one and can be deleted before any arithmetic. \emph{A hidden variable with no
observed descendants is irrelevant} --- a useful check to run before computing
anything.

\smallskip
\textbf{Step 2: the remaining sum.} With $r$ for \emph{Rushed}, $b$ for
\emph{Rollback}:
\[
  P(r \mid b) = \alpha\, P(r) \sum_{j} P(j) \sum_{w} P(w) \sum_{d}
  P(d \mid r,j,w)\, P(b \mid d)
\]

\textbf{Step 3: the inner sum.} For each combination of $r, j, w$ let
$q = P(\mathit{Defect} \mid r,j,w)$. Then
\[
  \sum_{d} P(d \mid r,j,w) P(b \mid d) = 0.70\,q + 0.05\,(1-q)
  = 0.05 + 0.65\,q
\]
so every row of the big table reduces to one number. With $r$ true the four $q$
values are $0.95, 0.80, 0.75, 0.50$, giving $0.6675, 0.5700, 0.5375, 0.3750$.

\smallskip
\textbf{Step 4: weight by the priors.} $P(j) = 0.40$, $P(w) = 0.25$:
\begin{align*}
  &0.40 \times 0.25 \times 0.6675 \;+\; 0.40 \times 0.75 \times 0.5700 \\
  &{}+\; 0.60 \times 0.25 \times 0.5375 \;+\; 0.60 \times 0.75 \times 0.3750 \\
  &= 0.06675 + 0.17100 + 0.08063 + 0.16875 \;=\; 0.48713
\end{align*}
Multiply by $P(r) = 0.30$: the unnormalised weight for $r$ true is $0.146138$.

\smallskip
\textbf{Step 5: the other branch.} With $r$ false the four $q$ values are
$0.70, 0.40, 0.30, 0.05$, giving $0.5050, 0.3100, 0.2450, 0.0825$; weighting by the
same priors gives $0.21738$, and multiplying by $P(\neg r) = 0.70$ gives
$0.152163$.

\smallskip
\textbf{Step 6: normalise.}
\[
  P(r \mid b) = \frac{0.146138}{0.146138 + 0.152163}
  = \frac{0.146138}{0.298301} = 0.4899
\]

\smallskip
\textbf{What the work shows.} Steps 3 and 4 were done once for $r$ true and once
for $r$ false, and the inner structure was identical both times. Variable
elimination is exactly the observation that this shared structure need not be
recomputed --- and on six variables the saving is not worth the bookkeeping, while
on \dsref{ve-scaling}'s chain it is the difference between possible and impossible.
\end{worked}

\begin{pitfall}
\begin{tight}
  \item \textbf{Reading arrows as ``causes'' and inference as following them.} The
        arrows point from cause to effect; the query here runs backwards, from
        rollback to specification. Both directions are ordinary conditional
        probability.
  \item \textbf{Forgetting that the missing arrows are the claims.} No edge between
        \emph{Rushed} and \emph{Junior} asserts independence. If rushed work is
        usually given to junior developers, the network is simply wrong.
  \item \textbf{Using rejection sampling with rare evidence.} $70\%$ of samples
        discarded here; with evidence of probability $0.01$ it is $99\%$. Use
        likelihood weighting.
  \item \textbf{Assuming variable elimination is always faster.} It lost on this
        network, $74$ operations against $32$. Its advantage is asymptotic and
        depends on the elimination order.
  \item \textbf{Expecting exact inference to scale in general.} It is linear on a
        chain and NP-hard on general graphs.
  \item \textbf{Summing over hidden variables that cannot matter.} A hidden variable
        with no observed descendants sums to one. Delete \emph{Hotfix} before
        computing, not after.
\end{tight}
\end{pitfall}

%---------------------------------------------------------------------
\begin{lab}[Diagnosing a Defect's Cause]
The network of \dsref{bayesnet}, four inference methods, the explaining-away
measurement, and the scaling comparison of \dsref{ve-scaling}.
\begin{lstlisting}[language=Python]
"""Chapter 15 lab -- inference in a Bayesian network.

Four ways to compute P(Rushed | Rollback): two exact, two by
sampling. Then explaining away, and the scaling comparison that
justifies variable elimination.
"""

import itertools
import random

# node -> (parents, table mapping parent values to P(node = True))
NET = {
    "Rushed":   ([], {(): 0.30}),
    "Junior":   ([], {(): 0.40}),
    "WeakTest": ([], {(): 0.25}),
    "Defect":   (["Rushed", "Junior", "WeakTest"], {
        (True, True, True): 0.95,   (True, True, False): 0.80,
        (True, False, True): 0.75,  (True, False, False): 0.50,
        (False, True, True): 0.70,  (False, True, False): 0.40,
        (False, False, True): 0.30, (False, False, False): 0.05,
    }),
    "Hotfix":   (["Defect"], {(True,): 0.90, (False,): 0.10}),
    "Rollback": (["Defect"], {(True,): 0.70, (False,): 0.05}),
}
ORDER = ["Rushed", "Junior", "WeakTest", "Defect", "Hotfix", "Rollback"]


def p_true(net, node, a):
    parents, table = net[node]
    return table[tuple(a[p] for p in parents)]


def prob(net, node, value, a):
    p = p_true(net, node, a)
    return p if value else 1 - p


# --- 1. exact: enumeration -----------------------------------------
def enumerate_query(net, order, query, evidence):
    hidden = [n for n in order if n != query and n not in evidence]
    totals, terms = {True: 0.0, False: 0.0}, 0
    for qv in (True, False):
        for combo in itertools.product([True, False],
                                       repeat=len(hidden)):
            a = dict(evidence)
            a[query] = qv
            a.update(dict(zip(hidden, combo)))
            pr = 1.0
            for n in order:
                pr *= prob(net, n, a[n], a)
            totals[qv] += pr
            terms += 1
    z = totals[True] + totals[False]
    return totals[True] / z, terms


# --- 2. exact: variable elimination --------------------------------
class Factor:
    def __init__(self, vars_, table):
        self.vars, self.table = list(vars_), table


def make_factor(net, node, evidence):
    allv = net[node][0] + [node]
    keep = [v for v in allv if v not in evidence]
    table = {}
    for combo in itertools.product([True, False], repeat=len(allv)):
        a = dict(zip(allv, combo))
        if any(v in evidence and a[v] != evidence[v] for v in allv):
            continue
        table[tuple(a[v] for v in keep)] = prob(net, node, a[node], a)
    return Factor(keep, table)


def multiply(f1, f2, ops):
    vars_ = f1.vars + [v for v in f2.vars if v not in f1.vars]
    table = {}
    for combo in itertools.product([True, False], repeat=len(vars_)):
        a = dict(zip(vars_, combo))
        table[combo] = (f1.table[tuple(a[v] for v in f1.vars)]
                        * f2.table[tuple(a[v] for v in f2.vars)])
        ops[0] += 1
    return Factor(vars_, table)


def sum_out(f, var, ops):
    i = f.vars.index(var)
    keep = [v for v in f.vars if v != var]
    table = {}
    for combo, val in f.table.items():
        k = combo[:i] + combo[i + 1:]
        table[k] = table.get(k, 0.0) + val
        ops[0] += 1
    return Factor(keep, table)


def var_elim(net, order, query, evidence):
    ops = [0]
    factors = [f for f in (make_factor(net, n, evidence) for n in order)
               if f.vars]
    for h in [n for n in order if n != query and n not in evidence]:
        involved = [f for f in factors if h in f.vars]
        if not involved:
            continue
        factors = [f for f in factors if h not in f.vars]
        prod = involved[0]
        for f in involved[1:]:
            prod = multiply(prod, f, ops)
        factors.append(sum_out(prod, h, ops))
    prod = factors[0]
    for f in factors[1:]:
        prod = multiply(prod, f, ops)
    i = prod.vars.index(query)
    totals = {True: 0.0, False: 0.0}
    for combo, val in prod.table.items():
        totals[combo[i]] += val
    z = totals[True] + totals[False]
    return totals[True] / z, ops[0]


# --- 3. rejection sampling -----------------------------------------
def rejection(net, order, query, evidence, n, rng):
    kept = hits = 0
    for _ in range(n):
        a = {}
        for node in order:
            a[node] = rng.random() < p_true(net, node, a)
        if all(a[k] == v for k, v in evidence.items()):
            kept += 1
            hits += a[query]
    return (hits / kept if kept else None), kept


# --- 4. likelihood weighting ---------------------------------------
def likelihood_weighting(net, order, query, evidence, n, rng):
    num = den = 0.0
    for _ in range(n):
        a, w = {}, 1.0
        for node in order:
            if node in evidence:
                a[node] = evidence[node]
                w *= prob(net, node, a[node], a)     # weight, not sample
            else:
                a[node] = rng.random() < p_true(net, node, a)
        den += w
        if a[query]:
            num += w
    return num / den, den


def make_chain(n):
    """X1 -> X2 -> ... -> Xn: a service dependency chain."""
    net = {"X1": ([], {(): 0.3})}
    for i in range(2, n + 1):
        net["X%d" % i] = (["X%d" % (i - 1)],
                          {(True,): 0.8, (False,): 0.1})
    return net, ["X%d" % i for i in range(1, n + 1)]


if __name__ == "__main__":
    print("variables: %d   full joint entries: %d"
          % (len(ORDER), 2 ** len(ORDER)))
    print("numbers the network stores: %d"
          % sum(len(NET[n][1]) for n in ORDER))
    print()

    ev = {"Rollback": True}
    exact, terms = enumerate_query(NET, ORDER, "Rushed", ev)
    ve, ops = var_elim(NET, ORDER, "Rushed", ev)
    print("P(Rushed)                = %.4f   (the prior)" % 0.30)
    print("P(Rushed | Rollback)     = %.5f" % exact)
    assert abs(exact - ve) < 1e-9
    print()
    print("method                          answer    work")
    print("-" * 48)
    print("enumeration (exact)           %9.5f  %4d terms" % (exact, terms))
    print("variable elimination (exact)  %9.5f  %4d ops" % (ve, ops))

    rng = random.Random(0)
    rej, kept = rejection(NET, ORDER, "Rushed", ev, 20000, rng)
    rng = random.Random(0)
    lw, _ = likelihood_weighting(NET, ORDER, "Rushed", ev, 20000, rng)
    print("rejection, 20000 samples      %9.4f  %4d kept (%.0f%%)"
          % (rej, kept, 100.0 * kept / 20000))
    print("likelihood weighting, 20000   %9.4f  all samples used" % lw)
    print()
    print("   errors: rejection %.4f, weighting %.4f"
          % (abs(rej - exact), abs(lw - exact)))
    assert abs(lw - exact) < 0.01 and abs(rej - exact) < 0.02

    # variable elimination LOSES on a network this small
    assert ops > terms
    print()
    print("Note that variable elimination used MORE work than")
    print("enumeration here (%d against %d). Its advantage is" % (ops, terms))
    print("asymptotic, and six variables is not asymptotic.")

    # --- explaining away -------------------------------------------
    print()
    print("=== explaining away ===")
    prior = 0.30
    given_d, _ = enumerate_query(NET, ORDER, "Rushed", {"Defect": True})
    given_dj, _ = enumerate_query(NET, ORDER, "Rushed",
                                  {"Defect": True, "Junior": True})
    print("   P(Rushed)                       = %.4f" % prior)
    print("   P(Rushed | Defect)              = %.4f   (rises)" % given_d)
    print("   P(Rushed | Defect, Junior)      = %.4f   (FALLS)"
          % given_dj)
    assert given_d > prior > 0 and given_dj < given_d
    print()
    print("   The junior developer already accounts for the defect,")
    print("   so less explaining is left for the rushed spec to do.")
    print("   Two causes of one effect are independent until the")
    print("   effect is seen, and negatively dependent afterwards.")
    print("   No production rule can express a conclusion that gets")
    print("   WEAKER when a condition is added.")

    # --- and where variable elimination earns its keep -------------
    print()
    print("=== scaling, on a chain of services ===")
    print("  n   enumeration   VE ops")
    print("-" * 30)
    for n in (6, 10, 14, 18):
        net, order = make_chain(n)
        e = {order[-1]: True}
        v, vops = var_elim(net, order, "X1", e)
        if n <= 14:
            ex, et = enumerate_query(net, order, "X1", e)
            assert abs(ex - v) < 1e-9
            print("%3d %13d %8d" % (n, et, vops))
        else:
            print("%3d %13s %8d" % (n, "infeasible", vops))
    print()
    print("Enumeration doubles with every service; variable")
    print("elimination grows linearly on a chain. At n=14 that is")
    print("8192 terms against 186 operations.")
\end{lstlisting}
Then add an arrow from \emph{Junior} to \emph{Rushed} --- suppose rushed work really
is given to junior developers --- and re-run. The explaining-away effect changes
size, because the two causes are no longer independent to begin with. That is the
modelling decision of the alertbox, made concrete.
\end{lab}

%---------------------------------------------------------------------
\begin{checkpoint}
\begin{tightnum}
  \item The full joint over these six variables needs $64$ numbers and the network
        stores $15$. State precisely what the network asserts in order to get away
        with the smaller number, and name one assertion in \dsref{bayesnet} that
        could be false in a real organisation.
  \item Learning that there is a defect raised belief in a rushed specification from
        $0.30$ to $0.5281$; additionally learning that a junior developer was
        involved \emph{lowered} it to $0.4304$. Explain why, and say why a
        forward-chaining rule base cannot reproduce this.
  \item Rejection sampling discarded $70\%$ of its samples here. Describe the
        circumstance in which it becomes useless, and say what likelihood weighting
        does differently.
\end{tightnum}
\end{checkpoint}

%---------------------------------------------------------------------
\begin{chaptersummary}
\begin{tight}
  \item A \term{Bayesian network} is a directed acyclic graph with a conditional
        probability table at each node. It asserts that each variable is
        conditionally independent of its non-descendants given its parents, which
        licenses the chain rule and lets $15$ numbers stand for a joint of $64$.
  \item \emph{The missing arrows are the claims.} They are substantive assertions
        about the domain and they can be wrong; drawing the graph is the modelling.
  \item \term{Enumeration} is exact and $O(2^{n})$. \term{Variable elimination} is
        exact and computes shared subexpressions once --- but it \emph{lost} on this
        six-node network, $74$ operations against $32$ terms. Its advantage is
        asymptotic: on a fourteen-service chain it is $186$ operations against
        $8{,}192$ terms.
  \item Exact inference in general Bayesian networks is NP-hard; the cost is
        governed by the largest intermediate factor.
  \item \term{Rejection sampling} discards every sample disagreeing with the
        evidence --- $70\%$ here, and nearly everything when the evidence is rare.
        \term{Likelihood weighting} fixes the evidence and weights instead, so every
        sample counts.
  \item \term{Explaining away}: two causes of a common effect are independent until
        the effect is observed and negatively dependent afterwards. Belief in a
        rushed specification rose to $0.5281$ on seeing a defect and fell to
        $0.4304$ on also learning of a junior developer.
  \item That behaviour is why one builds a network rather than a rule base: it
        requires a conclusion to become \emph{weaker} as a condition is added,
        which no production rule expresses and \chref{expert}'s specificity
        strategy would get backwards.
\end{tight}
\end{chaptersummary}

\begin{reviewq}
  \item \textbf{(MCQ)} A Bayesian network asserts that each variable is conditionally
        independent of: (a)~all other variables (b)~its non-descendants, given its
        parents (c)~its parents (d)~its descendants, given the evidence.
  \item \textbf{(MCQ)} Explaining away means that: (a)~evidence for one cause raises
        belief in another (b)~evidence for one cause lowers belief in another
        (c)~causes are always independent (d)~effects are independent of causes.
  \item \textbf{(MCQ)} Likelihood weighting differs from rejection sampling by:
        (a)~sampling more variables (b)~fixing the evidence and weighting each
        sample (c)~using exact arithmetic (d)~eliminating variables.
  \item \textbf{(MCQ)} A hidden variable with no observed descendants:
        (a)~must be summed over (b)~sums to one and can be deleted
        (c)~makes the query intractable (d)~must be given evidence.
  \item \textbf{(Short)} Explain what it means to say ``the missing arrows are the
        claims'', using one missing arrow from \dsref{bayesnet}.
  \item \textbf{(Short)} Variable elimination was slower than enumeration on this
        network and far faster on a chain of fourteen. Explain both results with
        one principle.
  \item \textbf{(Short)} Why can a production rule not express explaining away?
        State what such a rule would have to do.
  \item \textbf{(Applied)} Compute $P(\mathit{Defect} \mid \mathit{Hotfix})$ by hand
        from the tables of \dsref{bayesnet}, showing the marginalisation over
        $\mathit{Rushed}$, $\mathit{Junior}$ and $\mathit{WeakTest}$. Then say
        whether $\mathit{Rollback}$ needs to be summed over, and why.
\end{reviewq}
